% 全局排版：正文小四宋体，两字首行缩进；勿用空格模拟缩进。
\usepackage[margin=25mm,footskip=12mm]{geometry}
\setmainfont{TeX Gyre Termes}
\linespread{1.18}
\setlength{\parindent}{2\ccwd}
\setlength{\parskip}{0pt}
\pagestyle{plain}

% 一级标题用四号黑体，二级用小四黑体；节号为 1、1.1。
\ctexset{
  section = {format=\raggedright\heiti\zihao{4},name={},number=\arabic{section},
    aftername=\hspace{0.75em},beforeskip=1.1ex plus 0.3ex,afterskip=0.7ex,
    afterindent=true},
  subsection = {format=\raggedright\heiti\zihao{-4},name={},
    number=\arabic{section}.\arabic{subsection},aftername=\hspace{0.75em},
    beforeskip=0.9ex plus 0.2ex,afterskip=0.5ex,afterindent=true}
}
\setcounter{secnumdepth}{2}

% 首页为论文式标题区，不另设封面，不重置页码。
\makeatletter
\renewcommand{\@maketitle}{%
  \begin{center}
    {\heiti\zihao{3}\@title\par}
    \vspace{0.6em}
    {\zihao{-4}\@author\par}
    \vspace{0.2em}
    {\zihao{5}\@date\par}
  \end{center}%
}
\newcommand{\ReportPDFTitle}{\@title}
\newcommand{\ReportPDFAuthor}{\@author}
\makeatother

% 摘要与关键词采用五号宋体，标签用黑体。
\newenvironment{reportabstract}{%
  \par\begingroup\zihao{5}\noindent{\heiti 摘要：}\ignorespaces
}{\par\endgroup\smallskip}
\newcommand{\reportkeywords}[1]{\par\smallskip\noindent{\heiti 关键词：}#1\par}

\usepackage{amsmath,amssymb}
\usepackage{siunitx}
\usepackage{graphicx}
\usepackage{booktabs}
\usepackage{caption}
\DeclareCaptionFont{reportfive}{\zihao{5}}
\captionsetup{font=reportfive,labelsep=quad,justification=centering}
\captionsetup[figure]{position=bottom}
\captionsetup[table]{position=top}

% 顺序编码引文。保留 2015 样式以兼容较旧 TeX Live；不是 2025 版样式。
\usepackage[backend=biber,style=gb7714-2015]{biblatex}
\addbibresource{references.bib}
\renewcommand*{\bibfont}{\zihao{5}}
\usepackage[hidelinks,unicode]{hyperref}
\AtBeginDocument{\hypersetup{pdftitle={\ReportPDFTitle},pdfauthor={\ReportPDFAuthor}}}
